%这段需要修改 先放着
\begin{figure}[htbp]
    \centering
    \begin{tikzpicture}[
        node distance=1.5cm and 2.2cm,
        every node/.style={font=\small},
        box/.style={draw, rounded corners, minimum width=3.2cm, minimum height=1.1cm, align=center, fill=blue!6},
        arr/.style={-Stealth, thick, color=blue!60!black},
        dashedbox/.style={draw, dashed, rounded corners, inner sep=0.25cm, fit=#1, color=blue!40}
    ]
    
    % 第一排
    \node[box] (basic) {了解\\流动性风险\\[2pt]
    \textbf{交易/融资}\\
    杠杆\\
    抵押物
    };
    
    \node[box, right=of basic] (assetliab) {资产负债管理\\[2pt]
    \textbf{负债端/资产端}\\
    IS gap\\
    duration gap
    };
    
    \node[box, right=of assetliab] (liquidity) {流动性管理\\与国际业务\\[2pt]
    \textbf{储备/监督/定价}\\
    CIP迷失\\
    美元短缺
    };
    
    \node[box, right=of liquidity] (tools) {管理工具\\[2pt]
    \textbf{预警/日常/压力}\\
    流动性测试\\
    应急预案
    };
    
    % 箭头
    \draw[arr] (basic) -- (assetliab);
    \draw[arr] (assetliab) -- (liquidity);
    \draw[arr] (liquidity) -- (tools);
    
    % 虚线大框
    \node[dashedbox=(assetliab)(liquidity)(tools)] (bigbox) {};
    
    % 下方英文
    \node[below=1.2cm of bigbox, align=center, font=\small\itshape, text=blue!60!black] (eng) {Basic concept\\[2pt]
    Bank liquidity risk management\\(treasury department)};
    
    % 竖直箭头
    \draw[arr] (basic.south) -- ++(0,-0.7) -| (eng.north);
    
    \end{tikzpicture}
    \caption{银行流动性风险管理简明流程图}
    \end{figure}
    %这段表格需要修改 先放着







\subsection{Introduction to liquidity risk}
\begin{itemize}
\item Solvency vs. liquidity
    \begin{itemize}
    \item Solvency: assets more than liabilities (equity is positive).
    \item Liquidity: ability of a company to make cash payments as they become due.
    \end{itemize}
\item Two type of liquidity
    \begin{itemize}
    \item Transaction liquidity: an asset is easy to exchange for other assets.
    \item Funding liquidity: is the ability to finance assets continuously at an acceptable borrowing rate.
    \end{itemize}
\end{itemize}



\subsubsection{Sources of liquidity risk}
\begin{itemize}
\item \textbf{Transaction liquidity risk(a.k.a. Market liquidity risk or trading liquidity risk): }the risk of moving the price of an asset adversely in the act of buying or selling it.
    \begin{itemize}
    \item Transaction liquidity risk is low if assets can be liquidated quickly, cheaply, and without moving the price too much.
    \end{itemize}
\item \textbf{Funding liquidity risk (a.k.a. Balance sheet risk) :} the risk that creditors either withdraw credit or change the terms.
    \begin{itemize}
    \item Borrower's credit quality is or perceived to be deteriorating and borrower's financial conditions are deteriorating.
    \end{itemize}
\item Systemic risk: the risk of a general impairment of the financial system.
    \begin{itemize}
    \item In situations of severe financial stress, the ability of the financial system to allocate credit, support markets in financial assets, and even administer payments and settle financial transactions may be impaired.
    \end{itemize}
\item These types of liquidity risk interact.
    \begin{itemize}
    \item Funding liquidity risk increases transaction liquidity risk: If collateral requirements or the cost of financing increase, the trader may have to unwind securities.
    \item Transaction liquidity risk increases funding liquidity risk: If a leveraged market participant have illiquid assets on its books, its funding will be in greater jeopardy.
    \item Systemic risk increases transaction and funding liquidity risk.
    \end{itemize}
\end{itemize}



\subsubsection{Measuring market liquidity}
\begin{itemize}
\item Financial institutions can sell the asset at the bid to market maker and buy at the offer from market maker.
    \begin{itemize}
    \item The market maker is likely to increase the bid-offer spread above a certain size.
    \item The expected transactions cost is the half-spread or midto-bid spread.
    \end{itemize}
\item For asset, measure of market liquidity is its bid-offer spread
    \begin{itemize}
    \item Dollar amount: Offer price - Bid price.
    \item Proportional spread: $s = \dfrac{\text{Offer price} - \text{Bid price}}{\text{Mid market price}}$
    \end{itemize}
\item Example
    \begin{itemize}
    \item Suppose a dealer quoted the bid price and the offer price for a particular bond at 98 and 102 dollars, respectively. Calculate the dollar amount spread(p) and proportional spread(s).
    \end{itemize}
\item Answer:
    \begin{itemize}
    \item Offer price - Bid price = 102 - 98 = 4 dollars
    \item $ \text{Mid price} = \dfrac{\text{offer price} + \text{bid price}}{2} = \dfrac{102 + 98}{2} = 100\ \text{dollars} $
    \item $ S = \dfrac{\text{offer price} - \text{bid price}}{\text{Mid market price}}
= \dfrac{102 - 98}{\dfrac{102 + 98}{2}} = 4\% $
    \end{itemize}
\item For a book, measuring market liquidity is Cost of liquidation
    \begin{itemize}
    \item In normal market conditions:
        \begin{align*}
            \text{Cost of liquidation} 
            = \sum_{i=1}^{n} \frac{1}{2} s_i\, P_i\, Q_i = \sum_{i=1}^{n} \frac{1}{2} (\text{offer price}_i - \text{bid price}_i)\, Q_i
            \end{align*}
            \begin{center}
                $s_i$ : proportional spread \\
                $P_i$ : mid price \\
                $Q_i$ : liquidation volume
                \end{center}
    \end{itemize}
\item Example
    \begin{itemize}
    \item Tom's portfolio has two bond positions. The related information is shown in the table. What is the cost of liquidation of the whole portfolio in normal market?
        \begin{table}[htbp]
        \centering
        \begin{tabular}{|c|c|c|}
        \hline
        Portfolio & Bond A & Bond B \\
        \hline
        Market value(\$) & 40000 & 46000 \\
        \hline
        Proportional spread & 10\% & 20\% \\
        \hline
        \end{tabular}
        \end{table}
    \item Answer:
    \begin{align*}
        \text{Cost of liquidation of the whole portfolio} 
        &= \sum_{i=1}^{n} \dfrac{1}{2} s_i P_i Q_i \\
        &= \dfrac{1}{2} \times 10\% \times 40000 + \dfrac{1}{2} \times 20\% \times 46000 = 6600
        \end{align*}
 
    \end{itemize}

\item For a book, measuring market liquidity is Cost of liquidation
    \begin{itemize}
    \item In stressed market conditions:
        \begin{itemize}
        \item $\displaystyle
        \text{Cost of liquidation} = \sum_{i=1}^{n} \dfrac{1}{2} \left( \mu_i + \lambda \sigma_i \right) P_i Q_i
        $
        \begin{align*}
            \dfrac{1}{2} \left( \mu_i + \lambda \sigma_i \right) &: \text{Spread risk factor}
            \mu_i &: \text{mean of proportional spread} \\
            \sigma_i &: \text{standard deviation of proportional spread} \\
            \lambda &: \text{critical value under certain confidence level} \\
            P_i &: \text{mid price} \\
            Q_i &: \text{liquidation volume}
            \end{align*}
        \end{itemize}
    \end{itemize}
\item  Example
    \begin{itemize}
    \item  Given the following extra information, what is the 99\% confidential interval on the transaction cost per unit of bond A and the 99\% spread risk factor of bond B and what is the cost of liquidation of the whole portfolio in stressed market with confidence level of 99\%?
    \begin{table}[htbp]
        \centering
        \begin{tabular}{|c|c|c|}
        \hline
        Portfolio & Bond A & Bond B \\
        \hline
        Market value at mid price(\$) & 40000 & 46000 \\
        \hline
        Shares & 400 & 460 \\
        \hline
        Mean of spread($\mu_i$) & 10\% & 20\% \\
        \hline
        Standard deviation of spread($\sigma_i$) & 6\% & 8\% \\
        \hline
        \end{tabular}
        \end{table}
    \item    Answer
        \begin{itemize}
        \item 99\% expected transaction cost per unit of bond A:
        \[
        = \frac{1}{2} \times (10\% + 2.33 \times 6\%) \times (40000 \div 400) \times 1 = 11.99
        \]
        \item 99\% spread risk factor of bond B:
        \[
        = \frac{1}{2} \times (20\% + 2.33 \times 8\%) = 0.1932
        \]
        \item Cost of liquidation of the portfolio:
        \(
        = \sum_{i=1}^{n} \frac{1}{2} (\mu_i + \lambda \sigma_i) P_i Q_i
        \)
        \[
        = \frac{1}{2} \times (10\% + 2.33 \times 6\%) \times 40000 + \frac{1}{2} \times (20\% + 2.33 \times 8\%) \times 46000 = 13683.20
        \]
        \end{itemize}
    \end{itemize}
\item Liquidity-adjusted VaR is regular VaR plus the cost of liquidation.
    \begin{itemize}
    \item In normal market:
    \[
\text{Liquidity-adjusted VaR} = VaR + \sum_{i=1}^{n} \frac{1}{2} s_i P_i Q_i
\]
    \item In stressed market:
    \[
\text{Liquidity-adjusted VaR} = VaR + \sum_{i=1}^{n} \frac{1}{2} (\mu_i + \lambda \sigma_i) P_i Q_i
\]
    \end{itemize}
\item Example
    \begin{itemize}
    \item Allen holds a portfolio with 100,000 shares of a stock and the value of the position is \$4.0 million. The bid price of the stock is \$39 and the offer price is \$41. The stock's daily volatility is 2.59\%. Allen is interesting in calculating the liquidity-adjusted value at risk of the portfolio, Let's assume the stock's arithmetic returns are normally distributed (aka, normal VaR) and the expected daily return rounds to zero. Which is NEAREST to the position's one-day 99.0\% confident liquidityadjusted value at risk in normal market?
    \item Answer
        \begin{itemize}
        \item One day 99.0\% VaR:
        \(
        = \left| 0 - 2.59\% \times 2.326 \right| \times \$4.0\,\text{million} = \$240{,}973.60
        \)
        \item The liquidation cost:
        \(
        = \frac{1}{2} \times \dfrac{41.00 - 39.00}{40.00} \times \$4.0\,\text{million} = \$100{,}000.
        \)
        \item The one day 99.0\% LVaR:
        \(
        = 240{,}973.60 + 100{,}000 = \$340{,}973.60
        \)
        \end{itemize}
    
    \end{itemize}
\item VaR translation considering market liquidity:
    \begin{itemize}
    \item According to previews knowledge.
    \[
\text{T-day VaR} = \sqrt{T} \times VaR_{1\text{-}day}
\]
    \item Avoid adverse price impact, the portfolio position equally decreasing every day:
    \[
\text{Adjusted T-day VaR} = \sqrt{\frac{(1+T)(1+2T)}{6T}} \times VaR_{1\text{-}day}
\]
    \end{itemize}
\item Example
    \begin{itemize}
    \item Suppose a one-day VaR of the stock position is 23,000 dollars. To avoid adverse price impact, the fund manager plans to equally liquidate the position in 10 days. What is the adjusted 10-day VaR during the liquidation period?
    \item Answer:
    \begin{align*}
        \text{Adjusted 10-day VaR} 
        &= \sqrt{\frac{(1+T)(1+2T)}{6T}} \times VaR_{1\text{-}day} \\
        &= \sqrt{\frac{(1+10)(1+2\times10)}{6\times10}} \times 23{,}000 = 45{,}129.26\ \text{dollars}
        \end{align*}
    \end{itemize}

\end{itemize}



\subsubsection{Funding liquidity risk}
\begin{itemize}
    \item Funding liquidity risk is the risk of financial institution's ability to cannot finance its cash needs as they arise.
        \begin{itemize}
        \item The core function of a commercial bank is to take deposits and provide loans. In doing so, it carries out a liquidity and maturity, as well as credit transformation.
        \item Maturity mismatch leads to funding liquidity risk
            \begin{itemize}
            \item To manage funding liquidity risk, bank using asset liability management(ALM).
            \end{itemize}

        \end{itemize}
    \item Financial institutions that are solvent (i.e., have positive equity) and sometimes fail because of liquidity problems.
        \begin{itemize}
        \item Northern Rock
        \item Ashanti Goldfields
        \item Metallgesellschaft
        \end{itemize}
\end{itemize}



\subsubsection{Case of Northern Rock}
\begin{itemize}
    \item Case brief: Northern Rock was a fast-growing medium-sized mortgage bank based in the United Kingdom. The bank relied on selling short-term debt instruments for funding.
    
    Following the subprime crisis, the bank found it difficult to roll over the debt even though the bank was solvent. In November 2011, Northern Rock pie was bought from the British government for £747 million by the Virgin Group.
    \item Lessons: Liquidity can lead to serious problems even if the bank is solvent.
\end{itemize}



\subsubsection{Case of Ashanti Goldfields}
\begin{itemize}
    \item Case brief: Ashanti Goldfields had sought to protect its shareholders from gold price declines by selling gold forward. On September 26, 1999, 15 European central banks announced that they would limit their gold sales over the following 5 years. The price of gold jumped up over 25\%.
    
    Ashanti Goldfields was unable to meet margin calls, and this resulted in a major restructuring.
    \item Lessons: Liquidity problems can arise when hedging illiquid assets with contracts that are subject to margin requirements.
\end{itemize}



\subsubsection{Case of Metallgesellschaft}
\begin{itemize}
    \item Case brief: Metallgesellschaft sold a huge volume of 5 to 10 year heating oil and gasoline fixed-price supply contracts to its customers at six to eight cents above market prices. It hedged its exposure with long positions in futures. As it turned out, the price of oil fell and there were margin calls on the futures positions. The outcome was a loss to MG of \$1.33 billion.
    \item Lessons: Liquidity problems can arise when hedging illiquid assets with contracts that are subject to margin requirements
\end{itemize}


\subsection{liquidity challenges}


\subsubsection{Challenges faced by different institutions}
\begin{itemize}
\item Commercial banks:
    \begin{itemize}
    \item Bank run: due to fractional-reserve, if depositors are concerned about banks' liquidity, they will endeavor to redeem before other depositors and lenders.
    \end{itemize}
\item Security firms:
    \begin{itemize}
    \item Lenders abruptly withdrawing credit.
    \item Brokerage and clearing customers withdrawing deposits.
    \end{itemize}
\item Money Market Mutual Funds (MMMF):
    \begin{itemize}
    \item Sometimes MMMF had to liquidate assets with loss to meet huge redemption.
        \begin{itemize}
        \item Break the buck phenomenon.
        \end{itemize}
    \end{itemize}
\item Hedge funds:
    \begin{itemize}
    \item Hedge fund capital that can be redeemed does not only include which will be expected to experience large losses, but also those profitable or having low losses hedge funds.
        \begin{itemize}
        \item No access to wholesale funding and rely on other tools.
        \end{itemize}
    \end{itemize}
\end{itemize}



\subsubsection{Sources of Liquidity}
\begin{itemize}
    \item Holdings of cash and Treasury securities
    \item The ability to liquidate trading book positions
    \item The ability to borrow money at short notice
    \item The ability to offer favorable terms to attract retail and wholesale deposits at short notice
    \item The ability to securitize assets (such as loans) at short notice
    \item Borrowings from the central bank
\end{itemize}



\subsubsection{Regulations}
\begin{itemize}
\item Basel III introduced two liquidity risk requirements: the liquidity coverage ratio (LCR) and the net stable funding ratio (NSFR).
    \[
\text{\ LCR:}\quad
\frac{\textit{High quality liquid assets}}{\textit{Net cash outflows in a 30-day period}} \geq 100\%
\]

\[
\text{\ NSFR:}\quad
\frac{\textit{Available Amount of stable funding}}{\textit{Required amount of stable funding}} \geq 100\%
\]
\item BIS issued 17 principles for sound liquidity risk management.
    \begin{itemize}
    \item Identify the responsibilities, procedures, policies, strategies, information disclosure, supervision etc.
    \end{itemize}
\end{itemize}



\subsubsection{Liquidity black holes}
\begin{itemize}
\item Also called "crowded exit": a situation where liquidity has dried up in a particular market because everyone wants to do the same type of trade at the same time.
    \begin{itemize}
    \item During the sell-off, liquidity dries up
        \begin{itemize}
        \item When positive feedback traders dominate the trading, market prices are liable to be unstable and the market may become one-sided and illiquid.
        \end{itemize}
    \end{itemize}
\end{itemize}



\subsubsection{Positive and negative feedback traders}
\begin{itemize}
    \item Negative feedback traders: buy when prices fall and sell when prices rise.
    \item Positive feedback traders: sell when prices fall and buy when prices nse.
\end{itemize}



\subsubsection{Reasons of positive feedback trading}
\begin{itemize}
    \item Trend trading
    \item Stop-loss rules
    \item Dynamic hedging
    \item Creating options synthetically
    \item Margins
    \item Predatory trading
    \item Relative value fixed income trade
\end{itemize}






\subsection{Collateral and leverage}

\subsubsection{Collateral market}
\begin{itemize}
\item Collateral plays an important role in credit transactions by providing security for lenders and thus ensuring the availability of credit to borrowers.
\item Transactions in collateral market:
    \begin{itemize}
    \item Margin loans
    \item Repurchase agreements
    \item Securities lending
    \item Total return swaps
    \end{itemize}
\item Risks in collateral market transactions:
    \begin{itemize}
    \item Market risk: price fluctuation; interest rate fluctuation
    \item Credit risk: credit quality of the collateral deteriorates
    \item Counterparty risk: counterparty may default
    \end{itemize}
\end{itemize}



\subsubsection{Leverage}
\begin{itemize}
\item \textbf{Leverage ratio} $L = \dfrac{A}{E} = \dfrac{D+E}{E} = 1 + \dfrac{D}{E}$
    
\item \textbf{Leverage effect} can be seen by writing out the relationship between asset returns $r_A$, equity returns $r_E$, and the cost of debt $r_D$:
    \begin{itemize}

    \item $\bm{r_E = r_A + \dfrac{D}{E} \times (r_A - r_D)}$

    \item $\bm{r_E = L \times r_A - (L-1) \times r_D = r_D + L \times (r_A - r_D)}$
    \end{itemize}

\item \textbf{Effect of increasing leverage is} $\dfrac{\partial r_E}{\partial L} = r_A - r_D$
\item Example 1: Leverage and Required Returns
    \begin{itemize}
    \item A company sets its hurdle rate of ROE at least 15\%. The following table describes some financial data of the bank. What is the minimum leverage need to achieve its goal?
    \begin{table}[htbp]
        \centering
        \begin{tabular}{|l|c|}
        \hline
        Items & Amount \\
        \hline
        Return of asset & 12\% \\
        \hline
        Cost of debt & 8\% \\
        \hline
        \end{tabular}
        \end{table}
    \item Answer
        \begin{itemize}
            \item $r_E = r_D + L \times (r_A - r_D)$
            \item $15\% \leq 8\% + L \times (12\% - 8\%)$
            \item $L \geq 1.75$
        \end{itemize}
    \end{itemize}
\item Example 2: Leveraged Returns to Housing
    \begin{itemize}
    \item Suppose house value is expected to rise 10\% a year, neglecting the rental income and net of property maintenance costs. If the down payment of the house buyer is 20\%, what is the annual rate of return to the homeowner by 6\% loan rate?
    \item Answer
    \[
\frac{\text{house price appreciation} - (1 - \text{down payment}) \times \text{loan rate}}{\text{down payment}}
=
\frac{10\% - (1 - 20\%) \times 6\%}{20\%} = 26\%
\]
    \end{itemize}
\item Leverage in loans, options and derivatives
    \begin{itemize}
    \item An economic balance sheet is construct to capture the implicit or embedded leverage in loans, options and derivatives.
    \end{itemize}   
\end{itemize}



\subsubsection{Leverage: margin loans}
\begin{itemize}
\item Margin lending has a straightforward impact on leverage.
\item he haircut determines the amount of the loan that is made:
    \begin{itemize}
    \item Haircut of h percent: (1- h) is lent against a given market value of margin collateral.
    \item Leverage of h percent: The leverage on a position with a haircut of h percent is 1/h .
    \end{itemize}
\end{itemize}


\subsubsection{Leverage: short positions}
\begin{itemize}
\item Lengthen the balance sheet: since both the value of the borrowed short securities and the cash generated by their sale appear on the balance sheet.
    \begin{itemize}
    \item Increase leverage: both asset and liability increases the same amount with a constant equity.
        \[
L = \frac{A}{E} = 1 + \frac{D}{E}
\]
    \end{itemize}
\end{itemize}



\subsubsection{Leverage: derivatives}
\begin{itemize}
\item Derivative securities are a means to gain an economic exposure to some asset or risk factor without buying or selling it outright.
    \begin{itemize}
    \item Mimic cash position: If the derivative creates a synthetic long (short) position, the economic balance sheet entries mimic those of a cash long (short) position.
    \end{itemize}
    
\end{itemize}



\subsubsection{Leverage: structured credit}
\begin{itemize}
\item Seniority ranking: Bond tranche > Mezzanine tranche> Equity tranche.
    \begin{itemize}
    \item The equity tranches get the remaining of the asset pool.
    \end{itemize}
\item Practice
    \begin{itemize}
    \item On opening day, Master Fund LP has the following economic balance sheet:
    \begin{table}[htbp]
        \centering
        \begin{tabular}{|l|c|}
        \hline
        Cash & \$100 million \\
        \hline
        Financial assets & \$200 million \\
        \hline
        Debt & \$180 million \\
        \hline
        Equity & \$120 million \\
        \hline
        \end{tabular}
        \end{table}
        Assume Master Fund LP finances a long position in \$100
        million worth of a stock at the margin requirement of 30\%. It
        invests \$30million of its own funds and borrows \$70 million
        from a commercial bank. What is the change of firm's
        leverage ratio after trade?  
    \item Answer  
        \begin{itemize}
        \item 初始杠杆$\displaystyle = \frac{\textit{asset}}{\textit{equity}} = \frac{300}{120} = 2.5$。

        \item 交易后：
        
        \item 资产 $= 70\,(\textit{cash}) + 300\,(\textit{financial assets}) = \$370$
        
        \item 负债 $= 180\,(\textit{debt}) + 70\,(\textit{margin loan}) = \$250$
        
        \item 权益 $= \text{资产 - 负债} = 370 - 250 = 120$
        
        \item 交易后杠杆$\displaystyle = \dfrac{\textit{asset}}{\textit{equity}} = \dfrac{370}{120} = 3.08$
        
        \item 因此，杠杆从2.5变成3.08。
        \end{itemize}
    \end{itemize}
    
\end{itemize}


